\section{Related Work}

Representation learning in the space of NN weights has become a growing field recently. 
Several methods with different approaches to deal with weight spaces have been proposed to predict model properties such as accuracy \citep{unterthinerPredictingNeuralNetwork2020,eilertsenClassifyingClassifierDissecting2020, andreisSetbasedNeuralNetwork2023, zhangNeuralNetworksAre2023} or to learn the encoded concepts \citep{ashkenaziNeRNLearningNeural2022, deluigiDeepLearningImplicit2023}. 
%
Other work investigates the structure of trained weights on a fundamental level, using their eigen or singular value decompositions to identify training phases or predict properties \citep{martinTraditionalHeavyTailedSelf2019,MM20_SDM,martinPredictingTrendsQuality2021,martinImplicitSelfregularizationDeep2021, YTHx22_TR, mellerSingularValueRepresentation2023}.
Taking an optimization perspective, other work has investigated the uniqueness of the basis of trained NNs \citep{ainsworthGitReBasinMerging2022, brownPrivilegedConvergentBases2023}. Other work identifies subspaces of weights that are relevant, which motivates our work \citep{bentonLossSurfaceSimplexes2021, lucasMonotonicLinearInterpolation2021, wortsmanLearningNeuralNetwork2021, fortLargeScaleStructure2019}. The mode connectivity of trained models has been investigated to improve understanding of how to train models \citep{draxlerEssentiallyNoBarriers2018, nguyenConnectedSublevelSets2019, frankleLinearModeConnectivity2019}. 

A different line of work trains models to generate weights for target models, such as HyperNetworks \citep{haHyperNetworks2016, nguyenHyperVAEMinimumDescription2019, zhangGraphHyperNetworksNeural2019, knyazevParameterPredictionUnseen2021,knyazevCanWeScale2023b,kofinas2024graph}, with a recurrent backbone \citep{wangCompactOptimalDeep2023} as learned initialization \citep{dauphinMetaInitInitializingLearning2019} or for meta learning \citep{finnModelAgnosticMetaLearningFast2017,zhmoginovHyperTransformerModelGeneration2022,navaMetaLearningClassifierFree2022}.
While the last category uses data to get learning signals, another line of work learns representations of the weights directly. Hyper-Representations train an encoder-decoder architecture using reconstruction of the weights, with contrastive guidance, and has been proposed to predict model properties \citep{schurholtSelfSupervisedRepresentationLearning2021} or generate new models \citep{schurholtHyperRepresentationsPreTrainingTransfer2022, schurholtHyperRepresentationsGenerativeModels2022}. While previous work was limited to small models of fixed length, this paper proposes methods to decouple the representation learner size from the base model. Related approaches use convolutional auto-encoders \citep{berardiLearningSpaceDeep2022} or diffusion on the weights \citep{peeblesLearningLearnGenerative2022}.
